% Generated draft: requires mathematical and source review.
\documentclass[11pt]{article}
\usepackage[margin=25mm]{geometry}
\usepackage{fontspec}
\setmainfont{Latin Modern Roman}
\usepackage{amsmath,amssymb,mathtools,unicode-math}
\setmathfont{Latin Modern Math}
\usepackage{enumitem,needspace,xurl,hyperref}
\setlength{\parindent}{0pt}
\setlength{\parskip}{5pt}
\setlength{\emergencystretch}{3em}
\newcommand{\N}{\mathbb{N}}
\newcommand{\Z}{\mathbb{Z}}
\newcommand{\Q}{\mathbb{Q}}
\newcommand{\R}{\mathbb{R}}
\newcommand{\C}{\mathbb{C}}
\newcommand{\F}{\mathbb{F}}
\newcommand{\A}{\mathbb{A}}
\newcommand{\PP}{\mathbb{P}}
\newcommand{\E}{\mathbb{E}}
\newcommand{\eps}{\varepsilon}
\newcommand{\dd}{\,\mathrm{d}}
\begin{document}
\section*{UnsolvedMath ID: 800270}
\textbf{Source problem code:} catalogue entry
\subsection{Definitions and mathematical statement}
\textbf{Unresolved context.} The supplied statement does not define the measure \(\mu\), the underlying space, or the dynamical system. Consequently, it does not determine a standalone mathematical conjecture.

The following notation gives only a conditional template. If \(X\) is the intended topological space and \(\mu\) is a nonnegative Borel measure on \(X\), define
\[
\operatorname{supp}_X(\mu)=\{x\in X:\mu(U)>0\text{ for every open neighbourhood }U\text{ of }x\}.
\]
If \(f:X\to X\) is the intended dynamical map, a periodic point is a point \(p\) satisfying \(f^n(p)=p\) for some integer \(n\geq 1\); its periodic orbit is \(\{p,f(p),\ldots,f^{n-1}(p)\}\). Let \(\mathcal R_f\) denote the union of the repelling periodic orbits, with the meaning of ``repelling'' still requiring the omitted dynamical context.

The supplied question is then represented schematically by
\[
\operatorname{supp}_X(\mu)\stackrel{?}{=}\overline{\mathcal R_f}^{\,X}.
\]
The choices of \(X\), \(f\), and \(\mu\), their hypotheses and quantifiers, the intended closure topology, and the definition of repelling are not supplied. In particular, the nonnegative Borel measure convention above is conditional notation, not a recovered hypothesis of the source problem.
\subsection{Short English statement}
The question asks whether the support of an unspecified measure \(\mu\) equals the closure of the union of the repelling periodic orbits of an unspecified dynamical system. The omitted context is necessary for a precise statement. The registry labels this problem open; that label is supplied metadata, not an independently verified assessment.
\subsection{Sources}
\begin{itemize}
\item Local UnsolvedMath registry, supplied complete record: permanent ID \(800270\), source code \texttt{catalogue entry}, title \emph{Several Complex Variables Problem List --- Problem 69}.
\item Supplied source attribution: Buzzard, Fornaess, Krantz, \emph{The Several Complex Variables Problem List} (1995), Problem 69. Supplied source URL: \url{https://arxiv.org/abs/math/9509202}.
\item The record reports access on 2026-10-16 and extraction by \texttt{arxiv\_latex}. These provenance details are supplied metadata; the source has not been independently consulted or verified for this draft.
\end{itemize}
\end{document}
